\documentclass[10pt,a4paper]{amsart}
\usepackage[margin=19mm]{geometry}
\usepackage{amsmath,amssymb,mathtools}
\usepackage[scr=rsfs,cal=euler]{mathalfa}
\usepackage{xfrac}
\usepackage[unicode,hidelinks,bookmarks=false]{hyperref}
\newcommand{\ie}{i.e.\ }
\newcommand{\cf}{cf.\ }
\newcommand{\resp}{respectively\ }
\newcommand{\Subsection}{\subsection}
\providecommand{\nobreakdash}{\nobreak\hbox{-}}
\setlength{\emergencystretch}{2em}
\multlinegap0pt
\usepackage{bm}
\usepackage{bbm}
\usepackage{dsfont}
\usepackage{xspace}
\usepackage{cancel}
\usepackage{mathtools}
\usepackage{amsthm}
\makeatletter
\@ifundefined{Them}{\newtheorem{Them}{Theorem}}{}
\makeatother
\newcommand{\Bc}{\mathcal{B}}
\newcommand{\Dc}{\mathcal{D}}
\newcommand{\ZZ}{\mathbb{Z}}
\title[Characteristic quasi-polynomials of deletions of Shi arrange]{Characteristic quasi-polynomials of deletions of Shi arrangements of type C and type D}
\author{Akihiro Higashitani, Masato Konoike, Norihiro Nakashima, Satoshi Ono}
\date{Source: arXiv:2509.02043v2 (CC BY 4.0)}
\begin{document}
\maketitle

\noindent\emph{Source and licence.} Characteristic quasi-polynomials of deletions of Shi arrangements of type C and type D. Version arXiv:2509.02043v2 (CC BY 4.0), distributed under the Creative Commons Attribution 4.0 International licence, as recorded in the arXiv licence metadata for this exact version and shown on the abstract page https://arxiv.org/abs/2509.02043v2. Full rights notice: licenses/arxiv-2509.02043-CC-BY-4.0.txt.\par\smallskip

\noindent\textit{Editorial scope.} Counted unit: the Them whose source label is \texttt{thm-chara-quasi-TypeD} in the article (source0.tex lines 1058--1064, source label \texttt{thm-chara-quasi-TypeD}), delivered together with the article's own complete original proof of it (source0.tex lines 1065--1158, one \texttt{proof} environment of 94 lines). No mathematics is written, repaired or re-derived: statement and proof are the article's own text line for line. Required context retained uncounted: eq:typeB\_Shi (lines 128--130). Every definition, display and label of those lines is retained unchanged. Omitted from this package and not redistributed: 1--127, 131--1057, 1159--1802. Nothing inside a retained excerpt is omitted or altered: every retained body is delivered exactly as the article's lines are written, so no omission receipt of any kind accompanies this package. Citations: the retained excerpts carry no citation command at all, so no bibliography entry of the article is transcribed and no reference is redistributed. Reference adaptation: none was needed and none was made; every reference inside the delivered unit resolves inside it, so no unresolved reference or citation is produced. Printed slips kept literal: no printed word, symbol, hypothesis or number of the counted statement or of its proof was altered. Layout and class: the article's own class is \texttt{article} with packages including graphicx, amsmath, amssymb, bm, amsthm, mathrsfs, amscd, authblk, here, fancyhdr, tikz, color; the wrapper is the lane's pinned \texttt{amsart} editorial wrapper at 10pt on A4, which restates the theorem-like environments the retained text needs and adds the article's own macro definitions that the retained text needs (\texttt{\textbackslash Bc}, \texttt{\textbackslash Dc}, \texttt{\textbackslash ZZ}), with extra wrapper packages (bm, bbm, dsfont, xspace, cancel, mathtools). The delivered numbering is the unit's own. Assets: no retained span contains a figure environment, a graphics-inclusion command, a PDF-page inclusion, a table or a TikZ picture; no image file is copied into this package, the package declares no asset dependency and no image file is redistributed with it. Licence: version arXiv:2509.02043v2 of this article is distributed under the Creative Commons Attribution 4.0 International licence, as recorded in the arXiv licence metadata for this exact version and shown on the abstract page https://arxiv.org/abs/2509.02043v2. A full rights notice is delivered at licenses/arxiv-2509.02043-CC-BY-4.0.txt. Characters: every retained excerpt is pure ASCII, so the delivered engine, XeLaTeX with its default Latin Modern text font, reports no missing character.
\begin{align}\label{eq:typeB_Shi}
\left| M((\Bc_m)_q)\right| = (q-2m)^m. 
\end{align}

\begin{Them}\label{thm-chara-quasi-TypeD}
    The characteristic quasi-polynomial of the Shi arrangement of type D is 
\begin{align*}
\left|M(\Dc_q)\right|=(q-2m+2)^m
\end{align*}
for any $q\in\mathbb{Z}$ with $q\gg 0$.
\end{Them}

\begin{proof}
We construct a bijection between $M(\Dc_q)$ and $M(\Bc_{q+2})$. 
Once we establish it, we obtain the desired conclusion since we already know $\left|M(\Bc_q)\right|=(q-2m)^m$ by \eqref{eq:typeB_Shi}.  



\medskip

For each $\bm{x} \in M(\Dc_q)$ (resp. $\bm{y} \in M(\Bc_{q+2})$), since $\bm{x} \not\in H_{s,s',0}^-[q]$ (resp. $\bm{y} \not\in H_{s,s',0}^-[q+2]$) for any $s,s' \in [m]$ with $s \neq s'$, we have $x_s \neq x_{s'}$ (resp. $y_s \neq y_{s'}$). 
In particular, there is at most one $s \in [m]$ such that $x_s=0$ (resp. $y_s=q+1$). 
By taking this into account, we define the maps $\phi : M(\Dc_q) \to M(\Bc_{q+2})$ and $\psi : M(\Bc_{q+2}) \to M(\Bc_q)$ as follows:
\begin{itemize}
    \item For $\bm{x} \in M(\Dc_q)$, let
    \begin{align*}
        \phi(\bm{x})=
        \begin{cases}
            \bm{x} + \bm{1} + q\bm{e}_u &\text{if $x_u = 0$}, \\
            \bm{x} + \bm{1} &\text{otherwise}, 
        \end{cases}
    \end{align*}
    where $\bm{1} = (1, \ldots, 1)$, $\bm{e}_u$ denotes the $u$-th unit vector, and all resulting vectors are regarded as those in $\ZZ_{q+2}^m$. 
\end{itemize}
\begin{itemize}
    \item For $\bm{y} \in M(\Bc_{q+2})$, let
    \begin{align*}
        \psi(\bm{y})=
        \begin{cases}
            \bm{y} - \bm{1} - q\bm{e}_u &\text{if $y_u = q+1$}, \\
            \bm{y} - \bm{1} &\text{otherwise}, 
        \end{cases}
    \end{align*}
where all resulting points are regarded as vectors in $\ZZ_q^m$. 
\end{itemize}


We check that these maps are well-defined. Once we confirm it, since $\psi \circ \phi$ (resp. $\phi \circ \psi$) is the identity on $M(\Dc_q)$ (resp. $M(\Bc_{q+2})$), we can verify that these are bijective. 

\medskip

\noindent
For the well-definedness for $\phi$, let $\bm{x} \in M(\Dc_q)$.
It is clear that $\phi(\bm{x}) \in \ZZ_{q+2}^m$ by definition. 
%(We can also check for $\psi$ in the same way.) 
\begin{itemize}
%\item Assume $x_s \neq 0$ and $x_t \neq 0$. Then we have $(x_s + 1) - (x_t + 1) \neq 0$ from $\bm{x} \notin H_{s,t,c}^-[q]$, i.e. $x_s -x_t \neq 0$. 
%\item Assume that $x_s = 0$. (The case $x_t=0$ is similar.) Then we have $-x_t \neq - q$ from $\bm{x} \not\in H_{s,t,c}^-[q]$. Therefore, we obtain $(1 + q) - (x_t + 1) \neq 0$.
\item We see that $\phi(\bm{x}) \notin \bigcup_{s=1}^m (H_{s,0}[q+2] \cup H_{s,1}[q+2])$. 
    \begin{itemize}
        \item In the case where there is no index $u \in [m]$ such that $x_u = 0$, since $\bm{x} \in (\ZZ_q \setminus \{0\})^m$, we obtain $\phi(\bm{x}) = \bm{x} + \bm{1} \in (\ZZ_{q+2} \setminus \{0,1\})^m$. Thus, $\phi(\bm{x}) \notin \bigcup_{s=1}^m (H_{s,0}[q+2] \cup H_{s,1}[q+2])$ holds.
        \item In the case where there is $u \in [m]$ such that $x_u = 0$, the $u$-th entry of $\phi(\bm{x})$ is $q+1$, while each of the other entries is not $0$. Thus, $\phi(\bm{x}) \notin \bigcup_{s=1}^m (H_{s,0}[q+2] \cup H_{s,1}[q+2])$. 
    \end{itemize}
\item We see that $\phi(\bm{x}) \notin \bigcup_{1 \leq s<t \leq m}(H_{s,t,0}^+[q+2] \cup H_{s,t,1}^+[q+2])$.
    \begin{itemize}
        \item Assume $x_s \neq 0$ and $x_t \neq 0$. Since $x_s + x_t \in \ZZ_q \setminus \{0,1\}$, we obtain $(x_s + 1) + (x_t + 1) \in \ZZ_{q+2} \setminus \{0,1\}$. 
        \item Assume $x_s=0$ or $x_t=0$, say, $x_s=0$. Since $x_t \in \ZZ_q \setminus \{0,1\}$, we obtain $(q+1) + (x_t+1) \in \ZZ_{q+2}\setminus\{0,1\}$. 
    \end{itemize}
\item We see that $\phi(\bm{x}) \notin \bigcup_{1 \leq s<t \leq m}(H_{s,t,0}^-[q+2] \cup H_{s,t,1}^-[q+2])$.
    \begin{itemize}
        \item Assume $x_s \neq 0$ and $x_t \neq 0$. 
        Since $x_s - x_t \in \ZZ_q \setminus \{0,1\}$, we have $(x_s - 1) - (x_t - 1) \in \ZZ_{q+2} \setminus \{0,1\}$. 
        \item Assume that $x_s = 0$. (The case $x_t=0$ is similar.) 
        Then $-x_t \in \ZZ_q \setminus \{0,1\}$, i.e., $x_t \in \{1,\ldots,q-2\}$ as elements of $\ZZ$. 
        Thus, we obtain that $(q+1) - (x_t + 1) = q - x_t \in \{2,\ldots,q-1\}$ as $\ZZ$, i.e., $q-x_t \in \ZZ_{q+2} \setminus \{0,1\}$.  
    \end{itemize}
\end{itemize}

\medskip

\noindent
For the well-definedness for $\psi$, let $\bm{y} \in M(\Bc_{q+2})$.
Since $\bm{y} \in (\ZZ_{q+2} \setminus \{0,1\})^m$, we see that $\psi(\bm{y}) \in \ZZ_q^m$.   
\begin{itemize}
\item We see that $\psi(\bm{y}) \notin \bigcup_{1 \leq s<t \leq m}(H_{s,t,0}^+[q] \cup H_{s,t,1}^+[q])$.
    \begin{itemize}
        \item Assume $y_s \neq q+1$ and $y_t \neq q+1$. Since $y_s,y_t,y_s+y_t \in \ZZ_{q+2} \setminus \{0,1\}$, 
        we have $y_s + y_t \in \{4,5,\ldots,2q\} \setminus \{q+2,q+3\}$ as elements of $\ZZ$. 
        Thus, $y_s + y_t - 2 \in \{2,3,\ldots,2q-2\} \setminus \{q,q+1\}$, i.e., $(y_s-1)+(y_t-1)=y_s + y_t - 2 \in \ZZ_q \setminus \{0,1\}$. 
        \item Assume $y_s=q+1$ or $y_t=q+1$, say, $y_s=q + 1$. Then $y_t \in \ZZ_{q+2} \setminus \{0,1,2,q+1\}$ by $y_s+y_t \in \ZZ_{q+2} \setminus \{0,1\}, y_s - y_t \in \ZZ_{q+2} \setminus \{0\}$, and $y_t \in \ZZ_{q+2} \setminus \{0,1\}$. 
        Thus, $y_s + y_t \in \{q+4,q+5,\ldots,2q+1\}$ as $\ZZ$. Hence, we obtain $(y_s-q-1)+(y_t - 1)=y_t-1 \in \{2,\ldots,q-1\}$, 
        i.e., $y_t-1 \in \ZZ_q \setminus \{0,1\}$. 
    \end{itemize}
\item We see that $\psi(\bm{y}) \notin \bigcup_{1 \leq s<t \leq m}(H_{s,t,0}^-[q] \cup H_{s,t,1}^-[q])$.
    \begin{itemize}
        \item Assume $y_s \neq q+1$ and $y_t \neq q+1$. Since $y_s,y_t,y_s - y_t \in \ZZ_{q+2} \setminus \{0,1\}$, 
        we have $y_s-y_t \in \{-q+2,-q+3,\ldots,q-2\} \setminus \{0,1\}$ as $\ZZ$. 
        Thus, we obtain $(y_s -1) - (y_t - 1) =y_s-y_t \in \ZZ_q \setminus \{0,1\}$. 
        \item Assume that $y_s = q+1$. %(The case $y_t=q+1$ is similar.) 
        Then $y_t \in \ZZ_{q+2} \setminus \{0,1,2,q+1\}$ by $y_s - y_t \in \ZZ_{q+2} \setminus \{0\}, y_s + y_t \in \ZZ_{q+2} \setminus \{0,1\}$ and $y_t \in \ZZ_{q+2} \setminus \{0,1\}$. Thus, $q+1 - y_t \in \{2,3,\ldots,q-1\}$ as $\ZZ$. 
        Hence, we obtain $-(y_t - 1) \in \ZZ_q \setminus\{0,1\}$.
    \end{itemize}
\end{itemize}

These discussions imply the well-definedness of $\phi$ and $\psi$, as desired. 
\end{proof}

\end{document}
